\section{The oscillatory pulse and the patch obstruction}\label{sec:pulse}

We now prove \Cref{prop:patch} by a category argument. Assuming
density of a class of metrics admitting uniformly bounded heights,
we perturb a fixed metric inside a thin strip. Cap regularity and hyperbolic propagation give
smooth tangential data at both strip edges for every corresponding
height. A Taylor comparison for the fixed metric isolates the change
in the transverse derivative caused by the perturbation. Its
oscillatory moment is incompatible with smooth data at both edges.

\subsection{A bounded solution class}

The vector space of smooth symmetric tensors supported in
\(\mathcal P\) is a closed subspace of a smooth-function Fréchet
space. Its open subset \(\mathcal U\) is therefore a Baire space.
For \(M\in\N\), consider metrics admitting an admissible height with
\begin{equation}\label{eq:class-bounds}
\norm z_{C^8(S)}\le M,\qquad
M^{-1}\le\abs{H_{yy}}\le M,\qquad E\ge M^{-1}.
\end{equation}
Put
\begin{equation}\label{eq:speed-constants}
v_*=\sqrt{\kappa/2}\,M^{-3/2},\qquad
L=100(1+v_*^{-1}).
\end{equation}
For each rational \(q_0\), let \(\mathcal A_{M,q_0}\) be the
subclass in which such a height also satisfies
\begin{equation}\label{eq:q-class}
\abs{q(0,0)-q_0}<\frac1{100L}.
\end{equation}
These countably many classes cover every metric with an admissible
height. Indeed smoothness and the strict inequalities on the compact
square supply some \(M\), after which rational approximation supplies
\(q_0\). We will show that each class is nowhere dense.

Suppose, to the contrary, that the closure of one class contains a
nonempty open set \(\mathcal O\subset\mathcal U\). Fix
\(g_*\in\mathcal O\). All test metrics below lie in \(\mathcal O\)
and converge to \(g_*\) in \(C^8\). Metrics from the class can
be chosen arbitrarily close to each test metric in the smooth
topology. We may consequently use uniform low bounds for the metric,
its inverse, and the first derivatives of \(q\).
Since \(\abs q\le1/100\), a nonempty class satisfies
\(\abs{q_0}<1/20\).

Introduce linear coordinates
\begin{equation}\label{eq:shear}
\xi=x,\qquad \theta=y+q_0x.
\end{equation}
On a sufficiently small fixed neighborhood of the center, the
preceding low bounds imply
\begin{equation}\label{eq:central-bounds}
\abs{q-q_0}\le\frac1{10L},
\qquad \sqrt{-a}\ge v_*.
\end{equation}
The second inequality follows from \eqref{eq:negative-box}
and \eqref{eq:class-bounds}. In these coordinates,
\begin{equation}\label{eq:sheared-hessian}
\begin{gathered}
H_{\theta\theta}=H_{yy},\qquad
H_{\xi\theta}=(q-q_0)H_{yy},\\
H_{\xi\xi}=d_0H_{yy},\qquad d_0=(q-q_0)^2+a.
\end{gathered}
\end{equation}
Because \((10L)^{-1}\le v_*/1000\), \(d_0<0\) is separated
from zero. Solving the height equation in the other direction gives
\begin{equation}\label{eq:Q}
z_{\theta\theta}
=Q_g(\xi,\theta,z_i,z_{\xi\xi},z_{\theta\xi})
=\Gamma^i_{\theta\theta}z_i+
\frac{H_{\theta\xi}^2+(K\det g)(1-\abs{dz}_g^2)}
     {H_{\xi\xi}}.
\end{equation}
Its second-jet partial derivatives, at a solution, are
\begin{equation}\label{eq:hyperbolic-coefficients}
P_1=Q_{z_{\theta\xi}}=\frac{2(q-q_0)}{d_0},
\qquad S_1=Q_{z_{\xi\xi}}=-\frac1{d_0}>0.
\end{equation}
They have bounded first derivatives; \(S_1\) has a positive lower
bound. The following convenient speed bound follows from
\eqref{eq:central-bounds}:
\begin{equation}\label{eq:speed-bound}
\abs{P_1}+\sqrt{S_1}\le4/v_*<L/4.
\end{equation}

Choose \(r>0\) small enough that the slab
\begin{equation}\label{eq:slab}
-r\le\theta\le r,\qquad \abs\xi\le Lr
\end{equation}
and fixed margins lie in the preceding neighborhood and inside
\(\mathcal P\). We also require \(r<1/2\).
Along a flow line from \eqref{eq:flow},
\(D\theta=q_0-q\), whereas \(\theta=s\) when \(t=0\).
The global bound \(\abs q\le1/100\) keeps these flow segments
in that neighborhood after a further shrinking of \(r\). Thus throughout
the slab
\begin{equation}\label{eq:s-theta}
\abs{s-\theta}\le r/4.
\end{equation}
In fact the bound \(r/10\) follows from its \(t\)-length at most \(Lr\).

\subsection{The test metrics and smooth data at the pulse edges}

Choose a nonzero real \(\chi_0\in C_c^\infty((-Lr/8,Lr/8))\) and
a nonnegative \(\phi_0\in C_c^\infty((-1,1))\) with
\(\int\phi_0>0\). For small fixed \(\delta>0\) and integers
\(\tau\to\infty\), the test metrics are
\begin{equation}\label{eq:pulse}
\widetilde g_\tau
=g_*+\tau^{-N}\chi_0(\xi)\phi_0(\tau\theta/\delta)
                 \cos(\tau\xi)\dd\theta^2.
\end{equation}
Choose the fixed integer \(N>10\) larger than every derivative order
in one basic neighborhood of \(g_*\) contained in \(\mathcal O\).
For each fixed \(\delta\), the test metrics then belong to
\(\mathcal O\) for all sufficiently large \(\tau\). Only finitely
many seminorms are involved in this assertion.

By the closure assumption choose \(g_\tau\in\mathcal A_{M,q_0}\)
such that
\begin{equation}\label{eq:approximation}
\norm{g_\tau-\widetilde g_\tau}_{C^\tau(S)}
\le\tau^{-\tau},
\end{equation}
and choose a corresponding height \(z_\tau\) with
\eqref{eq:class-bounds} and \eqref{eq:q-class}.
We suppress its subscript when deriving estimates.
All comparisons using finitely many low derivatives hold uniformly
for large \(\tau\). Their limiting constants can be chosen independently
of the fixed small \(\delta\), since \(N>10\).

Away from the pulse rectangle
\[
\abs\theta\le\delta/\tau,\qquad \abs\xi\le Lr/8,
\]
the metrics have uniform bounds at every fixed derivative order.
By \eqref{eq:s-theta}, every point of that rectangle has
\(\abs s<r/2\) for large \(\tau\), in every solution's flow chart.
Consequently the image of every compact rectangle in
\((-2,2)\times(-2,-r/2)\), or in
\((-2,2)\times(r/2,2)\), misses the pulse support. On all such
images, \(g_\tau\) differs from the fixed metric \(g_*\) only
by the approximation error \eqref{eq:approximation}. This verifies
the metric hypothesis of \Cref{prop:cap} throughout both caps,
with \(b_*=-r/2\). The proposition gives bounds for every
derivative of \(z\) on the two cuts
\begin{equation}\label{eq:cap-cuts}
\theta=\pm r,\qquad \abs\xi\le Lr.
\end{equation}
These cuts lie in compact subsets of the respective flow-coordinate
caps, uniformly over the solutions.
\Cref{fig:pulse} shows the cuts and the pulse inside the
negative-curvature slab.

\begin{figure}[H]
\centering
\begin{tikzpicture}[font=\small,>=Latex,xscale=1.25,yscale=1.12]
  \fill[orange!8] (-1.6,-1.35) rectangle (1.6,1.35);
  \draw[gray] (-1.6,-1.35) rectangle (1.6,1.35);
  \draw[thick,blue!65!black] (-1.6,-1.35)--(1.6,-1.35);
  \draw[thick,blue!65!black] (-1.6,1.35)--(1.6,1.35);
  \draw[blue!65!black,dashed] (-1.6,-1.35)--(-.92,-.13);
  \draw[blue!65!black,dashed] (1.6,-1.35)--(.92,-.13);
  \draw[blue!65!black,dashed] (-1.6,1.35)--(-.92,.13);
  \draw[blue!65!black,dashed] (1.6,1.35)--(.92,.13);
  \fill[orange!45] (-.38,-.13) rectangle (.38,.13);
  \draw[thick] (-.38,-.13) rectangle (.38,.13);
  \draw[->,gray] (-1.8,0)--(1.88,0) node[right,black] {\(\xi\)};
  \draw[->,gray] (0,-1.70)--(0,1.75) node[above,black] {\(\theta\)};
  \node[anchor=east] at (-1.67,-1.35) {\(-r\)};
  \node[anchor=east] at (-1.67,1.35) {\(r\)};
  \node[fill=white,inner sep=1pt] at (0,1.05) {upper data};
  \node[fill=white,inner sep=1pt] at (0,-1.05) {lower data};
  \draw[->] (.9,.28)--(.34,.08);
  \node[anchor=west] at (.72,.35) {pulse};
\end{tikzpicture}
\caption{The hyperbolic slab, schematically and not to scale.
Cap regularity gives smooth data on the two cuts \(\theta=\pm r\).
The dashed lines indicate the shrinking spatial intervals used to
propagate those data toward both pulse edges. The pulse is supported
in \(\abs\theta<\delta/\tau\) and oscillates in \(\xi\).
The entire slab lies in negative curvature.}
\label{fig:pulse}
\end{figure}


\begin{lemma}[Propagation to the pulse]\label{lem:propagation}
At \(\theta=\pm\delta/\tau\), on \(\abs\xi\le Lr/2\),
the functions \(z,z_\theta\) have uniform bounds for every
tangential derivative.
\end{lemma}

\begin{proof}
We record the energy argument, which will also be used on the pulse.
If \(v\) satisfies
\[
v_{\theta\theta}-P_1v_{\theta\xi}-S_1v_{\xi\xi}
=b^\theta v_\theta+b^\xi v_\xi+R,
\]
with the above coefficient bounds and bounded \(b^i\), use the energy
density and flux
\[
e=\frac{v_\theta^2+S_1v_\xi^2}{2},
\qquad
j=\frac{P_1v_\theta^2}{2}+S_1v_\xi v_\theta.
\]
One has
\(\abs j\le(\abs{P_1}+\sqrt{S_1})e\).
Integrate on intervals whose endpoints move inward at speed \(L/4\).
Multiplication by \(v_\theta\) and integration by parts, with the
nonpositive boundary loss discarded, gives
\begin{equation}\label{eq:energy}
\mathcal E'(\theta)
\le C\mathcal E(\theta)
+\sqrt{2\mathcal E(\theta)}\,\norm R_{L^2},
\qquad \mathcal E(\theta)=\int e\dd\xi.
\end{equation}
Thus \(\sqrt{\mathcal E}\) is controlled by its initial value and
the time integral of \(\norm R_2\), with a Gronwall factor.
One can justify division at zero energy by first using
\(\sqrt{\mathcal E+\eta}\) and then letting \(\eta\downarrow0\).

Apply this estimate to \(v=\pd_\xi^jz\), \(j\ge8\), obtained by
differentiating \eqref{eq:Q}. The top second-jet terms give
the principal coefficients in \eqref{eq:hyperbolic-coefficients}.
If a second-jet argument receives \(j-1\) differentiations,
it gives \(v_\xi\) or \(v_\theta\); put these terms, together
with the direct first-jet linearization, into \(b^i\pd_iv\).
Their coefficients use only bounded low derivatives.
Put \(V=(z_\xi,z_\theta)\). Every remaining solution factor is
a spatial derivative of \(V\) of order at most \(j-1\),
already controlled by the induction. A factor of order \(j-1\)
requires at least \(j-2\) of the differentiations, since a
second-jet argument is already one spatial derivative of \(V\).
There cannot be two such factors, because \(2(j-2)>j\).
Bound that factor in \(L^2\), if it occurs. All other factors
have order at most \(j-2\), so the one-dimensional Sobolev
inequality and the available \(H^{j-1}\) bound for \(V\)
control them in supremum norm. This gives the \(L^2\) bound
for the remainder in \eqref{eq:energy}.
The same moving interval is used at every derivative order;
its length has a fixed positive lower bound, so the one-dimensional
Sobolev constants are uniform.

The initial first-jet bounds through spatial order seven follow
from the \(C^8\) class bound. For each larger order, the initial
energy on the cut \(\theta=-r\) is bounded by the cap estimate.
The metric derivatives are bounded on
\(-r\le\theta\le-\delta/\tau\), so \eqref{eq:energy}
closes the spatial induction there. The intervals start as
\([-Lr,Lr]\) and lose at most \(Lr/4\) at each end, leaving
room for Sobolev estimates on \(\abs\xi\le Lr/2\).
Reverse time and start at \(\theta=r\) to obtain the same bounds
at the upper pulse edge.
\end{proof}

\subsection{A comparison across the thin strip}

On \(\abs\xi\le Lr/2\), construct a smooth comparison \(z^0\)
on \(-\delta/\tau\le\theta\le\delta/\tau\) with the same data
\((z,z_\theta)\) at \(\theta=-\delta/\tau\), satisfying
\eqref{eq:Q} for the fixed metric \(g_*\) up to \(O(\tau^{-N})\).
For this purpose take the time Taylor polynomial of degree \(N+1\)
at the initial edge. Recursively determine its higher coefficients
from \eqref{eq:Q}, using \(g_*\). At step \(k\), for
\(0\le k\le N-1\), differentiating the right side \(k\) times
determines the time derivative of order \(k+2\) from those
through order \(k+1\) and their spatial derivatives. This is
possible because \(Q_g\) contains no second time derivative.
The denominator at the initial
data stays separated from zero by metric closeness and
\eqref{eq:sheared-hessian}. The coefficients and each fixed spatial
derivative are uniformly bounded by \Cref{lem:propagation}.
The denominator remains nonzero on the short strip. The recurrence
makes the residual and its first \(N-1\) time derivatives vanish
at the initial edge. Taylor's formula on a strip of length
\(2\delta/\tau\) gives the claimed \(O(\tau^{-N})\) bound
for fixed \(\delta\). In particular, \(z^0_\theta\) at the final
edge has bounds for all fixed spatial derivatives.

Put \(w=z-z^0\). Move \(Q_{g_*}(z)-Q_{g_*}(z^0)\) to the
linearized side, using mean partial derivatives along the segments
between their argument vectors
\((z_i,z_{\xi\xi},z_{\theta\xi})\). This yields
\begin{equation}\label{eq:w-equation}
w_{\theta\theta}-P_2w_{\theta\xi}-S_2w_{\xi\xi}
+\beta^i\pd_iw
=Q_{g_\tau}(z)-Q_{g_*}(z)+O(\tau^{-N}),
\end{equation}
with zero Cauchy data at the lower edge.
Here \(S_2>0\) has a uniform lower bound; the coefficients
are bounded and \(P_2,S_2\) have bounded first derivatives.
To see the positivity, the \(C^8\) bound keeps the actual argument
vector close throughout the strip to its initial value. The same is
true of the polynomial, by its coefficient bounds. The argument
vectors agree at the initial edge, because the Cauchy data agree.
At this common value, \(Q_{g_*,z_{\xi\xi}}>0\) by metric closeness and
\eqref{eq:hyperbolic-coefficients}. Continuity keeps it positive
on all the intervening segments, on which the denominator also
remains separated from zero. It is not necessary that those
segments consist of solutions.

A fixed speed bound larger than the one in \eqref{eq:speed-bound}
therefore applies to \eqref{eq:w-equation}. The limiting bounds
for the zeroth-order coefficient values and the positive lower
bound for \(S_2\) are independent of the fixed small \(\delta\).

\subsection{The curvature forcing and its moment}

In the linear coordinates \(\xi,\theta\), the curvature tensor formula
gives
\begin{equation}\label{eq:curvature-density}
K\det g
=\pd_{\xi\theta}g_{\xi\theta}
-\frac{\pd_\xi^2g_{\theta\theta}+\pd_\theta^2g_{\xi\xi}}2
+\mathcal L(g,\pd g),
\end{equation}
where \(\mathcal L\) is smooth for positive metric matrices.
Since \eqref{eq:pulse} changes only \(g_{\theta\theta}\),
its principal contribution is \(-\tfrac12\pd_\xi^2\)
of the scalar pulse. Consequently the right side of
\eqref{eq:w-equation} equals
\begin{equation}\label{eq:forcing}
\tau^{2-N}D_\tau\chi_0(\xi)\phi_0(\tau\theta/\delta)
                   \cos(\tau\xi)+o(\tau^{2-N}),
\qquad
D_\tau=\frac{1-\abs{dz}_{g_*}^2}{2H_{\xi\xi}(g_*,z)}.
\end{equation}
The error is uniform in space and time as \(\tau\to\infty\)
for each fixed \(\delta\). Metric changes involving at most one
derivative are \(O_\delta(\tau^{1-N})\); the approximation
\eqref{eq:approximation} and the Taylor residual are smaller.

There are constants \(c_D,C_D>0\), independent of fixed small
\(\delta\), such that for all large \(\tau\)
\[
c_D\le\abs{D_\tau}\le C_D,
\]
and \(D_\tau\) has one sign throughout the strip used.
Indeed the numerator is bounded below by the \(E\) bound,
uniform metric bounds, and metric closeness. The denominator
has the sign opposite to \(H_{yy}\), by \(d_0<0\), and stays
nonzero when \(g_\tau\) is replaced by \(g_*\).
The sign of \(H_{yy}\) is constant on the connected square.
Upper and lower bounds follow from the same low estimates.

Apply the zero-data estimate \eqref{eq:energy} to
\eqref{eq:w-equation}, shrinking intervals from
\([-Lr/2,Lr/2]\) at a fixed sufficiently large speed.
They contain \(\supp\chi_0\) throughout the strip for large
\(\tau\). Its source has \(L^2\) size \(O(\tau^{2-N})\),
so
\begin{equation}\label{eq:short-energy}
\norm{\pd w}_{L^2}\le(C_2+o(1))\delta\tau^{1-N}
\end{equation}
uniformly in strip time.
Here \(C_2\) can be chosen independently of small fixed
\(\delta\): the strip length is \(2\delta/\tau\), and every
Gronwall factor tends to one for fixed \(\delta\).
Finite coefficient derivative bounds depending on that \(\delta\)
are therefore harmless.

Integrate \eqref{eq:w-equation} against
\(\psi_\tau(\xi)=\chi_0(\xi)\cos(\tau\xi)\) over space and strip time.
The \(w_{\theta\theta}\) term is
\[
\int\psi_\tau(\xi)w_\theta(\xi,\delta/\tau)\dd\xi,
\]
because the initial data vanish. Both \(z_\theta\) and
\(z^0_\theta\) at the final edge have uniform spatial bounds of
every fixed order. Repeated integration by parts therefore makes
this term \(o(\tau^{1-N})\).

For each other second-order term, integrate one spatial derivative
by parts. The derivative of \(\psi_\tau\) costs \(O(\tau)\).
By \eqref{eq:short-energy} and the strip length, the total
absolute cost is at most
\begin{equation}\label{eq:moment-error}
(C_3+o(1))\delta^2\tau^{1-N}.
\end{equation}
The constant \(C_3\) is independent of small fixed \(\delta\);
terms with a derivative on a coefficient and the first-order terms
have an additional factor \(\tau^{-1}\) for fixed \(\delta\).
The spatial test is compactly supported, so there is no spatial
boundary contribution.

In contrast, \eqref{eq:forcing} paired with \(\psi_\tau\)
has absolute integral at least
\begin{equation}\label{eq:moment-lower}
(c_3+o(1))\delta\tau^{1-N},
\qquad c_3>0,
\end{equation}
with \(c_3\) independent of small fixed \(\delta\).
This follows from the constant sign and lower bound of \(D_\tau\),
the nonnegativity and positive integral of \(\phi_0\), and
\[
\int\chi_0(\xi)^2\cos^2(\tau\xi)\dd\xi
\longrightarrow \frac12\int\chi_0(\xi)^2\dd\xi>0.
\]
No convergence of \(D_\tau\) is needed.

Choose \(\delta>0\) so small that \(C_3\delta<c_3/2\),
and then let \(\tau\to\infty\). Equations
\eqref{eq:moment-error} and \eqref{eq:moment-lower} contradict
\eqref{eq:w-equation}.
The order of choices is summarized in \Cref{tab:pulse-choices}.
In particular, no convergence in every smooth seminorm was required
of the test metrics, and no convergence or common sign across
different \(\tau\) was required of \(D_\tau\).

\begin{table}[H]
\centering
\small
\renewcommand{\arraystretch}{1.18}
\begin{tabular}{@{}p{.15\textwidth}p{.36\textwidth}p{.43\textwidth}@{}}
\hline
Stage & Quantities fixed & Dependence permitted\\
\hline
First & Class, neighborhood, \(g_*\), slab, cutoffs, \(N\)
& Fixes leading constants \(C_2,C_3,c_3\), independent of
sufficiently small fixed \(\delta\).\\
Second & \(\delta>0\), with \(C_3\delta<c_3/2\)
& Subsidiary derivative bounds may depend on this fixed \(\delta\).\\
Last & Integers \(\tau\to\infty\)
& Thresholds and little-oh errors may depend on \(\delta\);
the leading constants do not.\\
\hline
\end{tabular}
\caption{Parameter order in the pulse contradiction. Every derivative
order has its own bound; no single bound for all orders is asserted.}
\label{tab:pulse-choices}
\end{table}

We have proved that every \(\mathcal A_{M,q_0}\) is nowhere dense.
Their countable union contains every metric admitting an admissible
height. Baire's theorem now proves \Cref{prop:patch}.
